\documentclass[11pt,a4paper]{article}
\usepackage[T1]{fontenc}
\usepackage[utf8]{inputenc}
\usepackage{lmodern}
\usepackage[margin=25mm]{geometry}
\usepackage{amsmath,amssymb,amsthm}
\usepackage{xcolor}
\usepackage[hidelinks]{hyperref}
\newtheorem{theorem}{Theorem}
\newtheorem{proposition}[theorem]{Proposition}
\newtheorem{lemma}[theorem]{Lemma}
\newtheorem{corollary}[theorem]{Corollary}
\theoremstyle{definition}
\newtheorem{definition}[theorem]{Definition}
\newtheorem{property}[theorem]{Property}
\theoremstyle{plain}
\newtheorem{remark}[theorem]{Remark}
\newsavebox{\statementbox}
\newenvironment{statementframe}{%
  \par\addvspace{0.55\baselineskip}\noindent
  \begingroup\setlength{\fboxrule}{0.3pt}\setlength{\fboxsep}{7pt}%
  \begin{lrbox}{\statementbox}%
  \begin{minipage}{\dimexpr\linewidth-2\fboxsep-2\fboxrule\relax}%
}{%
  \end{minipage}\end{lrbox}%
  {\color{black!45}\fbox{\color{black}\usebox{\statementbox}}}%
  \endgroup\par\addvspace{0.45\baselineskip}%
}
\DeclareMathOperator{\Tr}{Tr}
\DeclareMathOperator{\conv}{conv}
\DeclareMathOperator{\dist}{dist}
\newcommand{\C}{\mathbb C}
\newcommand{\Z}{\mathbb Z}
\newcommand{\F}{\mathbb F}
\newcommand{\E}{\mathcal E}
\newcommand{\hull}{\mathcal C}
\newcommand{\B}{\mathcal B}
\newcommand{\Wc}{\mathcal W}
\newcommand{\ket}[1]{\lvert#1\rangle}
\newcommand{\bra}[1]{\langle#1\rvert}
\newcommand{\leg}[2]{\left(\frac{#1}{#2}\right)}
\setlength{\emergencystretch}{2em}
\setlength{\parindent}{0pt}
\setlength{\parskip}{0.3em plus 0.1em minus 0.1em}
\title{Kirkwood--Dirac positive states of finite abelian Fourier pairs:\\
mixtures of pure states exhaust them only for cyclic groups\\ of prime-power order}
\author{Fable \& Astra\\[0.35em]
\normalsize Research led by Jean-Philippe Bruneton\\[-0.1em]
\small Universit\'e Paris Cit\'e, France}
\date{23 November 2026\\[0.2em]
\small First version, on $\Z_3\times\Z_3$: 17 November 2026}

\begin{document}
\maketitle
\begin{abstract}
Two incompatible quantum measurements have no joint probability distribution
in general, but the Kirkwood--Dirac (KD) representation assigns to every state
a table of complex quasiprobabilities with the correct marginals. When the two
measurements are the position basis and the Fourier basis attached to a finite
abelian group $G$, the states whose table is a genuine probability
distribution form a convex set, and its pure elements are known: they are
uniform superpositions over a coset of a subgroup, multiplied by a character.
This note determines when these pure states generate the whole set. Every
KD-positive state is a mixture of pure KD-positive states if and only if $G$
is cyclic of prime-power order. The ``if'' part is a theorem of De~Bi\`evre,
Langrenez and Radchenko, who also found counterexamples for $\Z_6$ and
$\Z_2\times\Z_2$; the ``only if'' part was open for all other groups. Our
proof has three ingredients. Subgroup and quotient embeddings preserve
exactly the trace-norm distance to the pure-state hull, which reduces the
question to the groups $\Z_p\times\Z_q$ with $p,q$ prime. On such a group the
four product bases form a Clauser--Horne--Shimony--Holt scenario, and an
explicit Bell-type witness is nonnegative on every pure KD-positive state and
negative on an explicit KD-positive state for all pairs $(p,q)$ outside a
list of eight. One of the eight, $\Z_2\times\Z_2$, is settled by hand; six are
settled by explicit states with exact rational KD tables whose positive
definiteness is certified in interval arithmetic; the last, $\Z_3\times\Z_3$,
by a state certified in rational arithmetic alone, at trace-norm distance at
least $111/20000$ from the hull.
\end{abstract}

\section{Introduction}
Two incompatible quantum measurements do not generally admit a joint
probability distribution. The Kirkwood--Dirac (KD) representation nevertheless
assigns a table of quasiprobabilities to a state and an ordered pair of
orthonormal bases. Its row and column sums are ordinary measurement
probabilities, but individual entries may be negative or complex. When every
entry is real and nonnegative, the table is a genuine probability distribution
and the state is called KD-positive for that pair of bases. This condition
depends on the chosen bases, and its failure is a resource: for multiqubit
systems, for instance, KD nonpositivity is necessary for a quantum
computational advantage \cite{Thio}. Among all quasiprobability
representations, the KD ones are singled out by the fact that they alone
produce natural conditional expectations \cite{SLBD}.

The pairs of bases studied here come from a finite abelian group $G$. The
first basis is labelled by the elements of $G$ and the second is its Fourier
transform, labelled by the characters of $G$. For a cyclic group this is the
discrete Fourier transform, and for $G=\Z_2^n$ it is the pair of
computational and Hadamard bases of $n$ qubits. For such a pair the pure
KD-positive states are completely known: each is the uniform superposition
over a coset of a subgroup, multiplied by a character \cite{DLR}. There are
finitely many of them, and every mixture of them is KD-positive.

The geometric question is whether the converse holds: is every KD-positive
mixed state a mixture of pure KD-positive states? An affirmative answer makes
the positive region a polytope whose vertices are listed by the pure-state
classification, and several arguments about KD-positive states need exactly
this description; the tight entropic uncertainty relations of
\cite{ThioEntropic}, for instance, are established under the hypothesis that
the KD-positive set is the hull of the basis projectors
\cite[Lemma~2]{ThioEntropic}. The answer depends on the basis pair. It is
affirmative for almost every pair of bases \cite{almost}, for the discrete
Fourier transform in prime dimension \cite{LAD} and in dimension $p^2$
\cite{Yang}, and for all cyclic groups of prime-power order \cite{DLR}. It is negative for some spin-one
bases \cite{LAD,convexroof} and for the Fourier pairs of $\Z_6$ and
$\Z_2\times\Z_2$ \cite{DLR}. De~Bi\`evre, Langrenez and Radchenko state the
case of an arbitrary finite abelian group as an open question
\cite[Section~8]{DLR}.

This note answers it: the cyclic groups of prime-power order are the only
ones for which the answer is affirmative. The obstruction is a Bell-type
inequality. Take $G=\Z_p\times\Z_q$ with $p$ and $q$ prime. The Hilbert space
is $\C^p\otimes\C^q$, and each factor carries two bases, position and Fourier.
A KD-positive state supplies a joint probability distribution for the four
resulting outcomes, so it obeys every classical correlation bound, such as
the Clauser--Horne--Shimony--Holt (CHSH) bound. Mixtures of pure KD-positive
states obey a strictly stronger bound, because in each pure KD-positive state
one of the two outcomes of each factor is uniformly distributed. Quantum
mechanics offers states whose KD table is still a probability distribution
but whose correlations violate the stronger bound; the room between the two
bounds exists because Tsirelson's bound $2\sqrt2$ exceeds the classical value
$2$. In the simplest case $G=\Z_2\times\Z_2$ the two bases of each qubit are
the eigenbases of the Pauli matrices $Z$ and $X$, and the state
\begin{equation}\label{eq:qubits}
 \rho=\frac14\Bigl(I-\frac13\bigl(Z\otimes Z+X\otimes Z+X\otimes X
 -Z\otimes X\bigr)\Bigr)
\end{equation}
is KD-positive without being a mixture of pure KD-positive states
(Proposition~\ref{prop:Z2Z2}). The general construction replaces $Z$ and $X$
by zero-mean observables that are diagonal in the two bases of $\Z_p$ and of
$\Z_q$, and what it measures is the gap between the two correlation bounds.

\paragraph{Setting.}
Let $G$ be a finite abelian group of cardinality $N$, with counting measure.
Its position basis is $(\ket g)_{g\in G}$ and its Fourier basis is
\[
 \ket\chi=N^{-1/2}\sum_{g\in G}\chi(g)\ket g,
 \qquad \chi\in\widehat G.
\]
For any operator $A$ define the KD array
\begin{equation}\label{eq:KD}
 Q_G(A)(g,\chi)=\bra\chi g\rangle\bra g A\ket\chi
 =\frac1N\sum_{h\in G}A_{g,h}\chi(h-g).
\end{equation}
The inverse and the Hilbert--Schmidt overlap identity are
\begin{align}
 A_{g,h}&=\sum_{\chi\in\widehat G}Q_G(A)(g,\chi)\chi(g-h),\label{eq:inverse}\\
 \Tr(A^*B)&=N\sum_{g,\chi}\overline{Q_G(A)(g,\chi)}Q_G(B)(g,\chi).
 \label{eq:overlap}
\end{align}
Write $\E_G$ for the density operators with nonnegative real KD array and
$\hull_G=\conv\{P\in\E_G:P\text{ is a rank-one projector}\}$.
We use the full trace norm $\|A\|_1=\Tr\sqrt{A^*A}$, without a factor $1/2$.

The classification of pure KD-positive states due to
De Bi\`evre, Langrenez and Radchenko \cite[Theorem~3.1]{DLR} is an input:
their vectors, up to global phase, are
\begin{equation}\label{eq:pures}
 \psi(g)=|H|^{-1/2}\gamma(g)\mathbf1_{a+H}(g),
 \qquad H\leq G,\quad a\in G,\quad\gamma\in\widehat G.
\end{equation}
With $H^\perp=\{\chi\in\widehat G:\chi|_H=1\}$, the projector $P$ onto $\psi$
has the array
\begin{equation}\label{eq:masks}
 Q_G(P)=\frac1N\,\mathbf1_{(a+H)\times(\gamma H^\perp)} ,
\end{equation}
supported on a ``rectangle'' with $|H|\,|H^\perp|=N$ points. By
\eqref{eq:overlap}, for a Hermitian operator $W$ with real array,
\begin{equation}\label{eq:average}
 \Tr(WP)=\text{the average of }N\,Q_G(W)\text{ over the rectangle }
 (a+H)\times(\gamma H^\perp).
\end{equation}

\paragraph{Results.}
The main statement is the following classification. The trivial group counts
as cyclic of prime-power order.

\begin{statementframe}
\begin{theorem}\label{thm:main}
Let $G$ be a finite abelian group. Then $\E_G=\hull_G$ if and only if $G$ is
cyclic of prime-power order.
\end{theorem}
\end{statementframe}

The ``if'' direction is \cite[Theorem~6.1]{DLR}. The ``only if'' direction
follows from three results proved below. Theorem~\ref{thm:lift} shows that the
zero extension from a subgroup and the pull-back from a quotient preserve the
trace-norm distance to the pure-state hull; a finite abelian group that is not
cyclic of prime-power order contains a subgroup isomorphic to
$\Z_p\times\Z_q$ with $p,q$ prime, so only these groups need a
counterexample. Theorem~\ref{thm:pq} provides one for every pair of primes,
by the CHSH mechanism described above. For $\Z_3\times\Z_3$ the counterexample
is the fully explicit state of Theorem~\ref{thm:example}, whose separation
from the hull is certified in rational arithmetic with the distance bound
$111/20000$.

\paragraph{Versions.}
The first version of this note (17 November 2026) contained the exact
$\Z_3\times\Z_3$ certificate, the transfer theorem and the resulting
classification for groups of order $2^a3^b$. The present version adds
Theorem~\ref{thm:pq} and, with it, the classification for all finite abelian
groups. The earlier material is kept unchanged in
Sections~\ref{sec:transfer} and~\ref{sec:z3}.

\section{Related work}
The first systematic study of mixed KD-positive states is due to Langrenez,
Arvidsson-Shukur and De~Bi\`evre. They proved that the KD-positive states are
exactly the mixtures of the basis projectors for the discrete Fourier
transform in prime dimension \cite[Theorem~1.1(iii)]{LAD}, and they exhibited
a pair of spin-one bases with KD-positive states outside the hull of the pure
ones \cite[Section~5.2]{LAD}. For a generic pair of bases the picture is the
simplest possible: with probability one the only pure KD-positive states are
the basis vectors, and every KD-positive state is a mixture of them
\cite[Theorem~1]{almost}. Mixed KD-positive states outside the pure-state
hull therefore require special basis pairs. For the spin-one pair, the convex
roof of the total KD nonpositivity equals one exactly on the pure-state hull,
so that it witnesses the separation faithfully
\cite[Section~4 and Appendix~B]{convexroof}.

The Fourier pairs are the natural structured case. For the discrete Fourier
transform in dimension $d$, Yang, Yao, Geng, Wang and Chen proved the
equality $\E=\hull$ for $d=p^2$ \cite[Theorem~1]{Yang} and conjectured it for
every $d$ \cite[Conjecture~1]{Yang}. For $d=pq$ with distinct primes they
showed that the operators with a real KD array are the real span of the pure
KD-positive projectors \cite[Lemma~4]{Yang}, and wrote that they believe the
equality but could not prove it. De~Bi\`evre, Langrenez and Radchenko moved
to arbitrary finite abelian groups: they classified the pure KD-positive
states \cite[Theorem~3.1]{DLR}, reduced the equality $\E_G=\hull_G$ to a
decomposition property of positive periodic functions
\cite[Theorem~5.2]{DLR}, and proved the equality for every cyclic group of
prime-power order \cite[Theorem~6.1]{DLR}. They refuted the conjecture of
\cite{Yang} by separated states for $\Z_6$ and for $\Z_2\times\Z_2$, obtained
by perturbing a state on a facet of the hull along its outward normal, with
positivity checked numerically \cite[Section~7]{DLR}, and left the general
abelian case open \cite[Section~8]{DLR}. Cai, Yang and Zheng gave a second
proof of the prime-power case based on a directed graph, proved
the equality on the real span of the pure states attached to any path of
that graph, and stated the structure for general $d$ as an open question
\cite[Theorems~1 and~2, Conclusion]{Cai}. Burkat and Strelchuk, classifying
the unitaries that preserve KD positivity, record that for $d=pq\neq6$ the
equality remained an open question to the best of their knowledge
\cite[Appendix~D]{Burkat}. Thio et al.\ found exotic states, KD-positive
and outside the pure-state hull, for the multiqubit groups $\Z_2^n$ with
$n\geq2$ by a related facet search \cite{Thio}. Spriet extended the pure-state
classification to second-countable locally compact abelian groups and
conjectured that the equality holds when the compact open subgroups are
totally ordered by inclusion, a condition that ``might also be necessary''
\cite[discussion after Eq.~(12)]{Spriet}. For finite groups this condition
singles out the cyclic groups of prime-power order.

Against this background, Theorem~\ref{thm:main} settles the finite case
completely. It refutes \cite[Conjecture~1]{Yang} for every $d=pq$ with
distinct primes, answers the open question of \cite[Section~8]{DLR}, and
proves that Spriet's condition is necessary for finite groups. The new cases
are the groups $\Z_p\times\Z_q$ with $pq\notin\{4,6\}$ and everything that
contains them; the subgroup reduction itself is elementary. The witness of
Section~\ref{sec:chsh} is of a different nature from the facet perturbations
used so far: it is given in closed form for all primes at once and identifies
the obstruction as a Bell inequality. The first version of this note treated
$\Z_3\times\Z_3$ by the facet construction of \cite[Section~7]{DLR} with an
exact certificate; that example is kept because it comes with an explicit
distance to the hull.

A search of the literature citing \cite{DLR} and of public repositories of
machine-generated proofs, made on 23 November 2026, found no earlier proof of
Theorem~\ref{thm:pq} or of Theorem~\ref{thm:main}. This statement concerns
bibliographic priority and is independent of the mathematical verification
described in Section~\ref{sec:verif}.

\section{Transfer between a group, its subgroups and its quotients}
\label{sec:transfer}
The following specialization of a channel-retraction argument carries a
counterexample from a subgroup or a quotient to the whole group, with its
distance to the hull. Its general norm-contraction mechanism is standard; the
point to verify is the action on the KD hulls.

\begin{statementframe}
\begin{theorem}\label{thm:lift}
Let $H\leq G$ and $K=G/H$, with $m=|H|$. Let $V:\ell^2(H)\to\ell^2(G)$
be zero extension and let $J:\ell^2(K)\to\ell^2(G)$ be the isometry
$(Jf)(g)=m^{-1/2}f(g+H)$. Each isometry sends KD-positive states to
KD-positive states and sends its pure-state hull into $\hull_G$.
Moreover, for $E=V$ with $T=H$, or $E=J$ with $T=K$, and every density
operator $\omega$ on $\ell^2(T)$,
\begin{equation}\label{eq:isometry}
 \dist_1(E\omega E^*,\hull_G)=\dist_1(\omega,\hull_T).
\end{equation}
\end{theorem}
\end{statementframe}

\begin{proof}
Equation \eqref{eq:KD} first gives the two formulas
\begin{align}
 Q_G(V\omega V^*)(g,\chi)
 &=\begin{cases}\dfrac mN Q_H(\omega)(g,\chi|_H),&g\in H,\\0,&g\notin H,
 \end{cases}\label{eq:subgroup}\\
 Q_G(J\omega J^*)(g,\chi)
 &=\begin{cases}\dfrac1m Q_K(\omega)(g+H,\bar\chi),&\chi|_H=1,\\0,&\chi|_H\ne1.
 \end{cases}\label{eq:quotient}
\end{align}
Here $\bar\chi$ is the induced character of $K$. These prove preservation
of KD positivity. Formula \eqref{eq:pures} proves preservation of pure
KD-positive projectors: for $V$, extend a character of $H$ to $G$; for
$J$, take the preimage in $G$ of the subgroup and coset in $K$.

We next check compression of a pure KD-positive vector $\psi$ of $G$.
For $V^*\psi$, its support is the intersection of a subgroup coset
$a+L$ with $H$. This is empty or a coset of $L\cap H$, so the compression
is zero or proportional to a vector in \eqref{eq:pures} for $H$.
For $J^*\psi$, sum the character on each intersection of a fibre of
$G\to K$ with $a+L$. Such an intersection is a coset of $H\cap L$.
If the character is nontrivial on $H\cap L$, the sum vanishes.
Otherwise its magnitude is constant on the coset $\pi(a+L)$, where
$\pi:G\to K$ is the quotient map, and its phase is a character of $\pi(L)$,
extendible to $K$. This compression is again zero or proportional to a pure
KD-positive vector of $K$.

Fix any pure $\tau\in\hull_T$. The completely positive trace-preserving map
\[
 \mathcal R_E(X)=E^*XE+\Tr[(I-EE^*)X]\tau
\]
satisfies $\mathcal R_E(\hull_G)\subseteq\hull_T$ by the compression check,
and $\mathcal R_E(E\omega E^*)=\omega$. The inclusion
$E\hull_TE^*\subseteq\hull_G$ and trace-norm invariance under an isometry
give the ``$\leq$'' direction of \eqref{eq:isometry}. For the reverse
direction, every $\sigma\in\hull_G$ satisfies
\[
 \|\omega-\mathcal R_E(\sigma)\|_1
 =\|\mathcal R_E(E\omega E^*-\sigma)\|_1
 \leq\|E\omega E^*-\sigma\|_1,
\]
by channel contraction on Hermitian operators. Taking infima finishes
the proof.
\end{proof}

\begin{corollary}\label{cor:reduction}
Let $T$ be a subgroup or a quotient of $G$, or a quotient of a subgroup. If
$\E_T\neq\hull_T$, then $\E_G\neq\hull_G$. Moreover, a finite abelian group
that is not cyclic of prime-power order contains a subgroup isomorphic to
$\Z_p\times\Z_q$ for some primes $p\leq q$.
\end{corollary}

\begin{proof}
The hull $\hull_T$ is a polytope, hence closed, so a state
$\omega\in\E_T\setminus\hull_T$ lies at positive distance from it, and
Theorem~\ref{thm:lift} gives a state of $\E_G$ at the same positive distance
from $\hull_G$. For the second assertion, if two distinct primes $p<q$ divide
$|G|$, Cauchy's theorem provides elements of orders $p$ and $q$, which
generate a subgroup isomorphic to $\Z_p\times\Z_q$. Otherwise $G$ is a
non-cyclic $p$-group, a direct sum of at least two nontrivial cyclic groups,
and elements of order $p$ in two summands generate $\Z_p\times\Z_p$.
\end{proof}

The sets $\E_G$ and $\hull_G$ are defined by the group structure alone: a
group isomorphism induces a unitary that maps position basis to position
basis and Fourier basis to Fourier basis. The group $\Z_p\times\Z_q$ with
$p\neq q$ may therefore be replaced by the cyclic group $\Z_{pq}$ whenever
convenient.

\section{A CHSH witness for $\Z_p\times\Z_q$}\label{sec:chsh}

\begin{statementframe}
\begin{theorem}\label{thm:pq}
For all primes $p\leq q$, the group $G=\Z_p\times\Z_q$ satisfies
$\E_G\neq\hull_G$.
\end{theorem}
\end{statementframe}

The case $p=q=3$ is Theorem~\ref{thm:example} of Section~\ref{sec:z3}. This
section proves all the others.

\subsection{The witness and the candidate state}
Here $N=pq$. Write a point of $G\times\widehat G$ as $x=(x_1,x_2,x_3,x_4)$,
where $(x_1,x_2)\in\Z_p\times\Z_q$ is the position and $(x_3,x_4)\in
\Z_p\times\Z_q$ labels the character
$\chi_{(x_3,x_4)}(g_1,g_2)=e^{2\pi i(x_3g_1/p+x_4g_2/q)}$. The Fourier basis
of $G$ is the tensor product of the Fourier bases $(\ket{f_\alpha})$ of
$\C^p$ and $(\ket{f_\beta})$ of $\C^q$. Let $u_1,u_3:\Z_p\to[-1,1]$ and
$u_2,u_4:\Z_q\to[-1,1]$ be functions with zero mean, and set
\begin{gather}
 A_0=\sum_{a\in\Z_p}u_1(a)\ket a\bra a,\qquad
 A_1=\sum_{\alpha\in\Z_p}u_3(\alpha)\ket{f_\alpha}\bra{f_\alpha},\notag\\
 B_0=\sum_{b\in\Z_q}u_2(b)\ket b\bra b,\qquad
 B_1=\sum_{\beta\in\Z_q}u_4(\beta)\ket{f_\beta}\bra{f_\beta},\notag\\
 \B=A_0\otimes B_0+A_1\otimes B_0+A_1\otimes B_1-A_0\otimes B_1,
 \qquad \Wc=I+\B,\qquad \rho_c=\frac{I-c\,\B}{N}\quad(c>0).
 \label{eq:chsh}
\end{gather}
The observables $A_0,A_1$ act on the first factor and $B_0,B_1$ on the
second; $\B$ is a CHSH operator, $\Wc$ is the witness and $\rho_c$ the
candidate state. Two families of functions are used. The \emph{alternating}
function is $v_2=(1,-1)$ for the prime $2$, and $v_p(0)=0$,
$v_p(x)=(-1)^x$ for $1\leq x\leq p-1$ when $p$ is odd. If $p\neq q$ we take
$u_1=u_3=v_p$ and $u_2=u_4=v_q$. If $p=q$ is odd we take the Legendre symbol
$u_1=u_2=u_3=\leg{\cdot}{p}$ and $u_4=-\leg{-1}{p}\leg{\cdot}{p}$. For a
function $u$ on $\Z_p$ put $\|u\|^2=\frac1p\sum_xu(x)^2$; this equals $1$ for
$v_2$ and $(p-1)/p$ for $v_p$ and for the Legendre symbol with $p$ odd. Let
\begin{equation}\label{eq:nu}
\begin{gathered}
 S(x)=u_1(x_1)u_2(x_2)+u_3(x_3)u_2(x_2)+u_3(x_3)u_4(x_4)-u_1(x_1)u_4(x_4),\\
 \nu=\bigl(\|u_1\|^2+\|u_3\|^2\bigr)\bigl(\|u_2\|^2+\|u_4\|^2\bigr).
\end{gathered}
\end{equation}

\begin{lemma}\label{lem:props}
With the notation above:
\begin{enumerate}
\item[(a)] $N\,Q_G(\Wc)=1+S$ and $N^2Q_G(\rho_c)=1-c\,S$, with $|S|\leq2$.
Hence $Q_G(\rho_c)\geq0$ whenever $c\leq\frac12$.
\item[(b)] If $(p,q)\neq(2,2)$, then $\Tr(\Wc P)\geq0$ for every pure
KD-positive projector $P$ of $G$.
\item[(c)] $\Tr(\Wc\rho_c)=1-c\,\nu$.
\end{enumerate}
\end{lemma}

\begin{proof}
(a) In $\bra g(A_i\otimes B_j)\ket\chi$, let an observable that is diagonal
in the position basis act on $\bra g$ and one that is diagonal in the Fourier
basis act on $\ket\chi$. This gives
$Q_G(A_0\otimes B_0)(x)=u_1(x_1)u_2(x_2)/N$,
$Q_G(A_1\otimes B_0)(x)=u_3(x_3)u_2(x_2)/N$, and similarly for the two other
terms, while $Q_G(I)=1/N$. Finally
$|S|=|u_2(u_1+u_3)+u_4(u_3-u_1)|\leq|u_1+u_3|+|u_3-u_1|\leq2$.

(b) By \eqref{eq:average}, $\Tr(\Wc P)$ is the average of $1+S$ over a
rectangle $(a+H)\times(\gamma H^\perp)$. For $H=0$, the position $(x_1,x_2)$
is fixed and $(x_3,x_4)$ is uniform on $\Z_p\times\Z_q$; the zero means
remove three of the four terms of $S$ and leave $1+u_1(x_1)u_2(x_2)\geq0$.
In the same way $H=G$, $H=\Z_p\times0$ and $H=0\times\Z_q$ give
$1+u_3u_4$, $1+u_3u_2$ and $1-u_1u_4$, all nonnegative. For $p\neq q$ there
are no other subgroups. For $p=q$ odd, the remaining subgroups are the lines
$L_s=\{(t,st)\}$ with $s\neq0$, for which $L_s^\perp=\{(x_3,x_4):x_3=-sx_4\}$.
On a rectangle attached to $L_s$ one has $x_2=sx_1+d$ with $x_1$ uniform and
$x_3=-sx_4+e$ with $x_4$ uniform and independent of $x_1$. The two mixed
terms of $S$ average to zero and
\[
 \Tr(\Wc P)=1+\mathbb E\bigl[u_1(x_1)u_2(sx_1+d)\bigr]
 +\mathbb E\bigl[u_3(-sx_4+e)u_4(x_4)\bigr].
\]
For $k\neq0$, the identity $\sum_x\leg{x(x+k)}{p}=-1$ holds, because for
$x\neq0$ one has $\leg{x(x+k)}{p}=\leg{1+k/x}{p}$ and $1+k/x$ runs over
$\F_p\setminus\{1\}$. Therefore $\frac1p\sum_x\leg xp\leg{sx+d}p$ equals
$\leg sp\frac{p-1}p$ if $d=0$ and $-\leg sp\frac1p$ if $d\neq0$. The first
expectation is thus $\leg sp\frac{p-1}p$ or $-\leg sp\frac1p$. By the choice
of $u_4$ and $\leg{-1}p\leg{-s}p=\leg sp$, the second expectation is
$-\leg sp\frac{p-1}p$ or $\leg sp\frac1p$. Whatever the sign of $\leg sp$,
one of the two expectations is at least $-\frac1p$ and the other at least
$-\frac{p-1}p$, so $\Tr(\Wc P)\geq0$.

(c) The zero means give $\Tr A_i=\Tr B_j=0$ and
$\Tr(A_0A_1)=\sum_{a,\alpha}u_1(a)u_3(\alpha)|\langle a|f_\alpha\rangle|^2
=\frac1p\sum_au_1(a)\sum_\alpha u_3(\alpha)=0$, and likewise $\Tr(B_0B_1)=0$.
Hence $\Tr\B=0$, the cross terms of $\Tr\B^2$ vanish, and
$\frac1N\Tr\B^2=\nu$. Then
$\Tr(\Wc\rho_c)=\frac1N\Tr\bigl((I+\B)(I-c\B)\bigr)=1-c\nu$.
\end{proof}

The operator $\rho_c$ is positive semidefinite when
$c\,\lambda_{\max}(\B)\leq1$, where $\lambda_{\max}$ denotes the largest
eigenvalue. Two bounds on $\lambda_{\max}(\B)$ are used.

\begin{lemma}\label{lem:tsirelson}
Let $A_0,A_1$ and $B_0,B_1$ be Hermitian contractions on two
finite-dimensional Hilbert spaces and let $\B$ be as in \eqref{eq:chsh}. Then
$\lambda_{\max}(\B)\leq2\sqrt2$. Moreover $\lambda_{\max}(\B)\leq2\kappa_A$,
where $\kappa_A=\max_\theta\|\cos\theta\,A_0+i\sin\theta\,A_1\|$ in operator
norm, and
symmetrically $\lambda_{\max}(\B)\leq2\kappa_B$ with
$\kappa_B=\max_\theta\|\cos\theta\,B_0+i\sin\theta\,B_1\|$.
\end{lemma}

\begin{proof}
Write $B_j$ for $I\otimes B_j$. For a unit vector $\psi$, the Cauchy--Schwarz
inequality gives
\[
 |\langle\psi,\B\psi\rangle|\leq\|(B_0-B_1)\psi\|+\|(B_0+B_1)\psi\|
 \leq\sqrt2\,\langle\psi,(2B_0^2+2B_1^2)\psi\rangle^{1/2}\leq2\sqrt2 ,
\]
which is Tsirelson's bound. For the second bound, dilate each $B_j$ to the
Hermitian unitary
\[
 R_j=\begin{pmatrix}B_j&\sqrt{I-B_j^2}\\\sqrt{I-B_j^2}&-B_j\end{pmatrix}
\]
on the doubled space. Both $R_0$ and $R_1$ compress to $B_0$ and $B_1$ under
the same isometry, so $\lambda_{\max}(\B)\leq\lambda_{\max}(\B')$, where
$\B'$ is built from $A_0,A_1,R_0,R_1$. Two Hermitian unitaries $R_0,R_1$
decompose their space into common invariant subspaces of dimension one or
two: if $w$ is an eigenvector of the unitary $R_0R_1$ with eigenvalue
$\lambda$, then $R_1w=\lambda R_0w$ and $R_1R_0w=\bar\lambda w$, so the span
of $w$ and $R_0w$ is invariant under both, and so is its orthogonal
complement. On a one-dimensional block $R_0=r_0$ and $R_1=r_1$ are signs and
$\B'=(r_0-r_1)A_0+(r_0+r_1)A_1$ is $\pm2A_0$ or $\pm2A_1$, of norm at most
$2\kappa_A$ (take $\theta=0$ or $\pi/2$). A two-dimensional block on which
$R_0$ or $R_1$ is $\pm I$ splits into one-dimensional blocks. On the others,
$R_0=\sigma_z$ and $R_1=\cos t\,\sigma_z+\sin t\,\sigma_x$ in a suitable
basis. Then $R_0-R_1=2\sin\frac t2\,\sigma_a$ and
$R_0+R_1=2\cos\frac t2\,\sigma_b$ for two anticommuting Pauli matrices
$\sigma_a,\sigma_b$, and in a basis where $\sigma_a=\sigma_x$ and
$\sigma_b=\pm\sigma_y$,
\[
 \B'=2\begin{pmatrix}0&M^*\\M&0\end{pmatrix},\qquad
 M=\sin\tfrac t2\,A_0\pm i\cos\tfrac t2\,A_1 ,
\]
whose largest eigenvalue is $2\|M\|\leq2\kappa_A$. Exchanging the roles of
the two factors gives the bound with $\kappa_B$.
\end{proof}

\begin{proposition}\label{prop:crit}
Let $(p,q)\neq(2,2)$, let $\Lambda\geq\lambda_{\max}(\B)$ and
$c=\min(\frac12,\frac1\Lambda)$. If $\nu>\max(2,\Lambda)$, then
$\rho_c\in\E_G\setminus\hull_G$.
\end{proposition}

\begin{proof}
The operator $\rho_c$ has trace one and is positive semidefinite since
$c\Lambda\leq1$; it is KD-positive by Lemma~\ref{lem:props}(a). By
Lemma~\ref{lem:props}(c), $\Tr(\Wc\rho_c)=1-c\nu<0$, whereas
$\Tr(\Wc\sigma)\geq0$ for every $\sigma\in\hull_G$ by
Lemma~\ref{lem:props}(b).
\end{proof}

\subsection{Case analysis}
Put $s_2=2$ and $s_p=2(p-1)/p$ for odd $p$, so that $\nu=s_ps_q$ for the
functions chosen above; $s_p$ increases with $p$ over the odd primes. Let
$\kappa_3$ and $\kappa_5$ be the constants $\kappa_A$ of
Lemma~\ref{lem:tsirelson} for the alternating functions $v_3$ and $v_5$. The
computation described in Section~\ref{sec:verif} certifies
$\kappa_3\leq1.25642$ and $\kappa_5\leq1.28772$.
\begin{center}
\renewcommand{\arraystretch}{1.25}
\begin{tabular}{llll}
\hline
pairs $(p,q)$ & lower bound on $\nu$ & $\Lambda$ & bound used\\
\hline
$(2,3)$ & $8/3\approx2.667$ & $2\kappa_3\leq2.5129$ & $\kappa$ bound on the $\Z_3$ factor\\
$(2,q)$, $q\geq5$ & $16/5=3.2$ & $2\sqrt2\approx2.8284$ & Tsirelson\\
$(3,q)$, $q\geq19$ & $48/19\approx2.5263$ & $2\kappa_3\leq2.5129$ & $\kappa$ bound on the $\Z_3$ factor\\
$(5,q)$, $q\geq7$ & $96/35\approx2.7429$ & $2\kappa_5\leq2.5755$ & $\kappa$ bound on the $\Z_5$ factor\\
$(p,q)$, $7\leq p<q$ & $240/77\approx3.1169$ & $2\sqrt2$ & Tsirelson\\
$(p,p)$, $p\geq7$ & $144/49\approx2.9388$ & $2\sqrt2$ & Tsirelson\\
\hline
\end{tabular}
\end{center}
In every row $\nu>\max(2,\Lambda)$, so Proposition~\ref{prop:crit} applies and
proves Theorem~\ref{thm:pq} for these pairs. The pairs left over are $(2,2)$,
$(3,3)$, $(5,5)$, $(3,5)$, $(3,7)$, $(3,11)$, $(3,13)$ and $(3,17)$.

\subsection{The group $\Z_2\times\Z_2$}
For two qubits the witness $\Wc$ takes the value $-1$ on one Bell state, and
a term $Y\otimes Y$ repairs it.

\begin{proposition}\label{prop:Z2Z2}
Let $G=\Z_2\times\Z_2$, so that $A_0=B_0=Z$ and $A_1=B_1=X$, and let
$\Wc'=I+\B-Y\otimes Y$. Then $\Tr(\Wc'P)\geq0$ on the twenty pure KD-positive
projectors, and for $\frac14<c\leq\frac{\sqrt2}4$ the operator $\rho_c$ is a
KD-positive state with $\Tr(\Wc'\rho_c)=1-4c<0$. The state \eqref{eq:qubits}
is $\rho_{1/3}$.
\end{proposition}

\begin{proof}
The subgroups $0$, $\Z_2\times0$, $0\times\Z_2$ and $G$ give sixteen product
states. On each of them exactly one of the four correlators of $\B$ equals
$\pm1$, the three others vanish, and $\langle Y\otimes Y\rangle=0$, so
$\Tr(\Wc'P)\in\{0,2\}$. The diagonal subgroup $\{00,11\}$ gives the four Bell
states, on which $(\langle Z\otimes Z\rangle,\langle X\otimes X\rangle,
\langle Y\otimes Y\rangle)$ is $(1,1,-1)$, $(1,-1,1)$, $(-1,1,1)$ or
$(-1,-1,-1)$ and the two mixed correlators vanish; the values of
$\Tr(\Wc'P)$ are $4,0,0,0$. The spectrum of $\B$ is $\{0,0,\pm2\sqrt2\}$, so
$\rho_c\geq0$ for $c\leq\sqrt2/4$, and $c\leq\frac12$ gives KD positivity by
Lemma~\ref{lem:props}(a). Finally $\Tr(\B\rho_c)=-c\nu=-4c$ and
$\Tr((Y\otimes Y)\rho_c)=0$.
\end{proof}

\subsection{Six pairs settled by explicit states}\label{sec:six}
For $(3,5)$, $(3,7)$, $(3,11)$, $(3,13)$, $(3,17)$ and $(5,5)$ the witness
$\Wc$ is kept and $\rho_c$ is replaced by a state of the form
\begin{align*}
 \rho&=\sum_{a,b}w_{12}(a,b)\,\ket a\bra a\otimes\ket b\bra b
 +\sum_{\alpha,b}w_{32}(\alpha,b)\,\ket{f_\alpha}\bra{f_\alpha}\otimes\ket b\bra b\\
 &\quad+\sum_{\alpha,\beta}w_{34}(\alpha,\beta)\,\ket{f_\alpha}\bra{f_\alpha}\otimes\ket{f_\beta}\bra{f_\beta}
 +\sum_{a,\beta}w_{14}(a,\beta)\,\ket a\bra a\otimes\ket{f_\beta}\bra{f_\beta},
\end{align*}
a real combination of the product pure KD-positive projectors, with rational
weights of denominator $10^7$ that sum to one and may be negative. By
\eqref{eq:masks},
\[
 N\,Q_G(\rho)(x)=w_{12}(x_1,x_2)+w_{32}(x_3,x_2)+w_{34}(x_3,x_4)+w_{14}(x_1,x_4),
\]
and $\Tr(\Wc\rho)$ is the corresponding weighted sum of the values
$1\pm u_iu_j$ of Lemma~\ref{lem:props}(b). Both are exact rational numbers.
The weights were obtained by semidefinite programming, mixed with the
maximally mixed state to gain a positivity margin, and rounded; the proof
uses only the stored weights. The checker verifies in exact arithmetic that
$Q_G(\rho)\geq0$, that $\Tr(\Wc\rho)<0$, and again that $\Tr(\Wc P)\geq0$ on
every pure KD-positive projector, and certifies $\rho>0$ in interval (ball)
arithmetic.
\begin{center}
\renewcommand{\arraystretch}{1.25}
\begin{tabular}{llll}
\hline
$(p,q)$ & $\Tr(\Wc\rho)$ & $\min N\,Q_G(\rho)$ & certified lower bound on $\lambda_{\min}(\rho)$\\
\hline
$(3,5)$ & $-246111/2500000$ & $333/2500000$ & $1.33\times10^{-4}$\\
$(3,7)$ & $-663653/5000000$ & $951/10000000$ & $9.52\times10^{-5}$\\
$(3,11)$ & $-334063/2000000$ & $151/2500000$ & $6.05\times10^{-5}$\\
$(3,13)$ & $-35569/200000$ & $4/78125$ & $5.12\times10^{-5}$\\
$(3,17)$ & $-943767/5000000$ & $391/10000000$ & $3.91\times10^{-5}$\\
$(5,5)$ & $-2181631/10000000$ & $799/10000000$ & $7.99\times10^{-5}$\\
\hline
\end{tabular}
\end{center}
Together with Proposition~\ref{prop:crit}, the case table,
Proposition~\ref{prop:Z2Z2} and Theorem~\ref{thm:example} below, this proves
Theorem~\ref{thm:pq}.

\section{An exact certificate for $\Z_3\times\Z_3$}\label{sec:z3}
For $G=\Z_3\times\Z_3$ the conjugate basis is obtained by applying the
three-dimensional Fourier transform to each qutrit, namely $F_3\otimes F_3$;
this is different from the Fourier transform $F_9$ of the cyclic group
$\Z_9$, for which $\E=\hull$. Each of the six subgroups of $G$ gives nine
pure KD-positive states, for a total of $54$, which include both product and
entangled states. The following statement gives a separated state by the
facet construction of \cite[Section~7]{DLR}, with a certificate that uses
rational arithmetic alone and yields a distance to the hull. In this section
$W$ denotes the witness of the theorem, normalized so that the hull lies
below the value $18$.

\begin{statementframe}
\begin{theorem}\label{thm:example}
For the position/Fourier pair of $G=\Z_3\times\Z_3$, there exist an explicitly
specified density operator $\rho>0$ and a Hermitian operator $W$ such that
\begin{align}
 \min Q_G(\rho)&=\frac{2291}{1935000}>0,\label{eq:minQ}\\
 \max_{\sigma\in\hull_G}\Tr(W\sigma)&=18,
 &\Tr(W\rho)&=\frac{206973}{10000}\approx20.6973,\label{eq:gap}\\
 \dist_1(\rho,\hull_G)&\geq\frac{111}{20000}.\label{eq:distance}
\end{align}
In particular, $\hull_G\subsetneq\E_G$.
\end{theorem}
\end{statementframe}

Identify $G$ with $\F_3^2$ and put $\zeta=e^{2\pi i/3}$. Rows and columns
are both ordered lexicographically as
$(0,0),(0,1),(0,2),(1,0),\ldots,(2,2)$; column $c$ represents the character
$\chi_c(g)=\zeta^{c\cdot g}$.
Define $\rho$ and $W$ by \eqref{eq:inverse} from
\begin{equation}\label{eq:Rdef}
 Q_G(\rho)=\frac{R}{7740000},\qquad Q_G(W)=M,
\end{equation}
where
\begin{equation}\label{eq:R}
\begingroup\setlength{\arraycolsep}{2.8pt}
R=\begin{pmatrix}
99613&123096&99613&99613&123096&99613&123096&56130&103096\\
99613&99613&123096&99613&99613&123096&123096&103096&56130\\
123096&123096&123096&123096&123096&123096&96130&76130&76130\\
99613&99613&123096&123096&76130&123096&123096&79613&79613\\
99613&99613&123096&123096&99613&99613&123096&56130&103096\\
103096&56130&79613&56130&79613&103096&76130&9164&32647\\
123096&99613&99613&99613&123096&99613&123096&103096&56130\\
123096&123096&76130&99613&123096&99613&123096&79613&79613\\
56130&103096&79613&103096&79613&56130&76130&32647&9164
\end{pmatrix},
\endgroup
\end{equation}
and
\begin{equation}\label{eq:M}
M=\begin{pmatrix}
-1&8&-1&-1&8&-1&8&-10&8\\
-1&-1&8&-1&-1&8&8&8&-10\\
8&8&8&8&8&8&-10&-10&-10\\
-1&-1&8&8&-10&8&8&-1&-1\\
-1&-1&8&8&-1&-1&8&-10&8\\
8&-10&-1&-10&-1&8&-10&-28&-19\\
8&-1&-1&-1&8&-1&8&8&-10\\
8&8&-10&-1&8&-1&8&-1&-1\\
-10&8&-1&8&-1&-10&-10&-19&-28
\end{pmatrix}.
\end{equation}

\subsection{Hermiticity, normalization and positivity}
All entries of $R$ are positive and their sum is $7740000$.
Equation \eqref{eq:inverse} yields $\rho=A+B\zeta$ with rational matrices
$A,B$. Direct substitution, using $\zeta^2=-1-\zeta$, gives
\[
 A_{ji}=A_{ij}-B_{ij},\qquad B_{ji}=-B_{ij},
\]
so $\rho$ is Hermitian. The same check applies to $W$. Also $\Tr\rho=1$,
and the smallest entry of $R$ is $9164$, proving \eqref{eq:minQ}.

Here is a short rational certificate for $\rho>0$. For $i\ne j$ set
\[
 u_{ij}=|A_{ij}-B_{ij}/2|+\frac78|B_{ij}|.
\]
Since $\rho_{ij}=(A_{ij}-B_{ij}/2)+i\sqrt3 B_{ij}/2$ and
$\sqrt3/2<7/8$, one has $|\rho_{ij}|\leq u_{ij}$.
For $s=(1,1,1,1,1,2,1,1,2)$, exact rational arithmetic gives
\begin{equation}\label{eq:margins}
\begin{split}
 m_i&=s_i\rho_{ii}-\sum_{j\ne i}s_j u_{ij},\\
 (m_i)_{i=1}^9&=\frac1{61920000}
 (1360135,1360135,1376216,940135,940135,\\[-2pt]
 &\hspace{36mm}330636,786094,786094,536024).
\end{split}
\end{equation}
Every entry is strictly positive. With $D=\operatorname{diag}(s)$, the
Hermitian matrix $D\rho D$ has positive diagonal and is strictly diagonally
dominant, hence is positive definite. Equivalently, Gershgorin applied to
$D^{-1}\rho D$ gives the useful explicit bound
\begin{equation}\label{eq:eiglower}
 \lambda_{\min}(\rho)\geq\min_i\frac{m_i}{s_i}
 =\frac{27553}{10320000}>0.
\end{equation}
The much larger numerical eigenvalue estimate is not used in this proof.

\subsection{Separation from all pure KD-positive mixtures}
The subgroups of $\F_3^2$ are the zero subgroup, four lines, and the full
group. For each subgroup $H$ the distinct pure projectors in
\eqref{eq:pures} have the arrays \eqref{eq:masks}, with $N=9$.
There are $|G/H||\widehat G/H^\perp|=9$ of them for each $H$,
and hence $54$ in total. Each rectangle has nine occupied entries.
By \eqref{eq:overlap}, $\Tr(WP)$ is the sum of the entries of $M$ on
that rectangle. Enumerating the six subgroups and their cosets gives the
multiset
\begin{equation}\label{eq:scores}
 \{\Tr(WP):P\text{ pure KD-positive}\}
 =\{18\ (43\text{ times}),\ -63\ (10\text{ times}),\ -144\ (1\text{ time})\}.
\end{equation}
This is a finite integer calculation, not a numerical optimization.
Linearity proves that the maximum over $\hull_G$ is $18$.
On the other hand, direct substitution gives
\[
 \Tr(W\rho)=\frac9{7740000}\sum_{g,c}R_{g,c}M_{g,c}
 =\frac{206973}{10000}\approx20.6973>18.
\]
This proves \eqref{eq:gap} and excludes $\rho$ from $\hull_G$.

For completeness the trace-norm bound has a conservative rational
certificate as well. Write $W=A^{(W)}+B^{(W)}\zeta$ by
\eqref{eq:inverse}. Then
\[
 \|W-18I\|_\infty
 \leq\max_i\sum_j
 \bigl(|A^{(W)}_{ij}-18\delta_{ij}|+|B^{(W)}_{ij}|\bigr)=486.
\]
Here Hermiticity allows the maximum absolute row sum to bound the
operator norm. For any $\sigma\in\hull_G$, the trace of $\rho-\sigma$
vanishes and H\"older's inequality yields
\[
 486\|\rho-\sigma\|_1
 \geq\Tr[(W-18I)(\rho-\sigma)]
 \geq\frac{26973}{10000}.
\]
Since $26973/(486\cdot10000)=111/20000$, this completes the proof of
Theorem~\ref{thm:example}.

\begin{remark}
The arrays also admit a geometric construction. If $\rho_0$ is the uniform
average of the $43$ pure projectors with score $18$ in \eqref{eq:scores},
then $Q_G(\rho)=Q_G(\rho_0)+M/20000$. This is the facet-perturbation
construction of \cite[Section~7.2]{DLR}; the witness was found numerically,
and the positivity proof above is independent of that search.
\end{remark}

\section{Proof of Theorem~\ref{thm:main} and consequences}
\begin{proof}[Proof of Theorem~\ref{thm:main}]
If $G$ is cyclic of prime-power order, $\E_G=\hull_G$ by
\cite[Theorem~6.1]{DLR}. Otherwise $G$ contains a subgroup isomorphic to
$\Z_p\times\Z_q$ with $p\leq q$ prime (Corollary~\ref{cor:reduction}), for
which the equality fails by Theorem~\ref{thm:pq}, and the failure passes to
$G$ by Corollary~\ref{cor:reduction}.
\end{proof}

Theorem~\ref{thm:lift} also transports distances, which gives a quantitative
statement wherever an explicit distance is available.

\begin{statementframe}
\begin{corollary}
Every finite abelian group with a $\Z_3^2$ subquotient admits a state in
$\E_G$ at trace-norm distance at least $111/20000$ from $\hull_G$.
In particular this holds for all noncyclic finite abelian $3$-groups.
\end{corollary}
\end{statementframe}
\begin{proof}
Compose the subgroup and quotient isometries in Theorem~\ref{thm:lift}.
A noncyclic abelian $3$-group contains a subgroup $\Z_3^2$.
\end{proof}

\begin{remark}
For groups of order $2^a3^b$ the classification already follows from
Theorem~\ref{thm:lift}, Theorem~\ref{thm:example} and the two examples of
\cite[Section~7]{DLR}: a noncyclic group has a noncyclic Sylow subgroup,
which contains $\Z_2^2$ or $\Z_3^2$, and a cyclic group with $a,b>0$ contains
$\Z_6$. This was the content of the first version of this note; it relied on
the numerical positivity checks of \cite{DLR} for $\Z_6$ and $\Z_2^2$, which
Proposition~\ref{prop:Z2Z2} and the first row of the case table now replace
by proofs.
\end{remark}

\begin{remark}
The state $\rho_c$ has a joint probability distribution for the four product
bases, so it satisfies the classical CHSH bound
$\Tr(\B\rho_c)\geq-2$, while mixtures of pure KD-positive states satisfy
$\Tr(\B\sigma)\geq-1$ for $(p,q)\neq(2,2)$. In the decomposition criterion of
\cite[Theorem~5.2]{DLR}, the obstruction comes from the four-cycle formed by
the subgroups $0$, $\Z_p\times0$, $G$ and $0\times\Z_q$, which is absent when
the subgroups form a chain, as they do for $\Z_{p^k}$. The analogous question
for infinite locally compact abelian groups, including the sufficiency
conjecture of \cite{Spriet}, is not addressed here.
\end{remark}

\section{Verification materials}\label{sec:verif}
The proofs of Theorems~\ref{thm:lift} and~\ref{thm:example},
Lemmas~\ref{lem:props} and~\ref{lem:tsirelson} and
Propositions~\ref{prop:crit} and~\ref{prop:Z2Z2} are written out above. Three
steps of Theorem~\ref{thm:pq} are computer-assisted: the bounds on $\kappa_3$
and $\kappa_5$, the six states of Section~\ref{sec:six}, and, as a
convenience, the arithmetic of Theorem~\ref{thm:example}. The public bundle
contains the corresponding checkers and data.

For $\Z_3\times\Z_3$, standard-library Python programs reconstruct all $54$
pure arrays, verify strict positivity and the witness value, and check the
diagonal-dominance certificate in exact arithmetic: the
\href{https://github.com/jpbruneton/open_quantum_problems/blob/main/research/entries/kd-z3x3/checks/verify_diagonal_certificate.py}{weighted diagonal-dominance checker},
the
\href{https://github.com/jpbruneton/open_quantum_problems/blob/main/research/entries/kd-z3x3/checks/verify_z3_certificate.py}{independent cyclotomic checker},
and the
\href{https://github.com/jpbruneton/open_quantum_problems/blob/main/research/entries/kd-z3x3/checks/verify_paper_data.py}{manuscript-data checker},
together with the
\href{https://github.com/jpbruneton/open_quantum_problems/blob/main/research/entries/kd-z3x3/data/z3x3-simple-certificate.json}{exact certificate}.

For the CHSH witness, the
\href{https://github.com/jpbruneton/open_quantum_problems/blob/main/research/entries/kd-z3x3/checks/certify_chsh_states.py}{state checker}
reads the stored rational weights of the six states
(\href{https://github.com/jpbruneton/open_quantum_problems/tree/main/research/entries/kd-z3x3/data}{files \texttt{chsh-state-p-q.json}})
and verifies, in exact rational arithmetic, the trace, the nonnegativity of
the KD array, the sign of $\Tr(\Wc\rho)$ and the nonnegativity of $\Wc$ on
every pure KD-positive projector. Positive definiteness is certified with Arb
ball arithmetic at 160 bits: the matrix is conjugated by a floating-point
approximation $U$ of its eigenvectors, and $U^*\rho U$ is shown to have a
positive diagonal and to be strictly diagonally dominant, which implies
$\rho>0$ whatever the quality of the approximation. The
\href{https://github.com/jpbruneton/open_quantum_problems/blob/main/research/entries/kd-z3x3/checks/certify_kappa.py}{$\kappa$ checker}
uses $\|\cos\theta\,A_0+i\sin\theta\,A_1\|^2=\lambda_{\max}\bigl(C_0+\cos2\theta\,
C_1+\sin2\theta\,C_2\bigr)$ with $C_0=\frac12(A_0^2+A_1^2)$,
$C_1=\frac12(A_0^2-A_1^2)$ and $C_2=\frac i2[A_0,A_1]$; it certifies an upper
bound at each of $4000$ equally spaced angles by the same ball-arithmetic
test and adds the Lipschitz term $\frac{\pi}{4000}(\|C_1\|_F+\|C_2\|_F)$,
where $\|\cdot\|_F$ is the Frobenius norm. The
\href{https://github.com/jpbruneton/open_quantum_problems/blob/main/research/entries/kd-z3x3/checks/check_chsh_analytic.py}{analytic cross-checker}
verifies Proposition~\ref{prop:Z2Z2} symbolically, the rational inequalities
of the case table, and Lemma~\ref{lem:props}(b) exhaustively on thirteen
pairs up to $(13,13)$. The common routines are in
\href{https://github.com/jpbruneton/open_quantum_problems/blob/main/research/entries/kd-z3x3/checks/chsh_common.py}{\texttt{chsh\_common.py}}.
These files complement the proofs given above.

\begingroup\small
\begin{thebibliography}{99}
\setlength{\itemsep}{2pt}
\bibitem{DLR}
S. De Bi\`evre, C. Langrenez and D. Radchenko,
\emph{The Kirkwood--Dirac representation associated to the Fourier transform
for finite abelian groups: positivity},
Ann. Henri Poincar\'e \textbf{27}, 3093--3112 (2026);
\href{https://doi.org/10.1007/s00023-025-01614-7}{doi:10.1007/s00023-025-01614-7};
\href{https://arxiv.org/abs/2501.12252v2}{arXiv:2501.12252v2}.

\bibitem{LAD}
C. Langrenez, D. R. M. Arvidsson-Shukur and S. De Bi\`evre,
\emph{Characterizing the geometry of the Kirkwood--Dirac-positive states},
J. Math. Phys. \textbf{65}, 072201 (2024);
\href{https://doi.org/10.1063/5.0164672}{doi:10.1063/5.0164672};
\href{https://arxiv.org/abs/2306.00086v1}{arXiv:2306.00086v1}.

\bibitem{Yang}
Y.-H. Yang, S. Yao, S.-J. Geng, X.-L. Wang and P.-Y. Chen,
\emph{Geometry of Kirkwood--Dirac classical states: a case study based on
discrete Fourier transform},
J. Phys. A: Math. Theor. \textbf{57}, 435303 (2024);
\href{https://doi.org/10.1088/1751-8121/ad819a}{doi:10.1088/1751-8121/ad819a};
\href{https://arxiv.org/abs/2404.09399v2}{arXiv:2404.09399v2}.

\bibitem{almost}
C. Langrenez, W. Salmon, S. De Bi\`evre, J. J. Thio, C. K. Long and
D. R. M. Arvidsson-Shukur,
\emph{Almost sure minimality of the set of experiments with classical
Kirkwood--Dirac representations},
Phys. Rev. A \textbf{113}, 062215 (2026);
\href{https://doi.org/10.1103/v7z4-qsz8}{doi:10.1103/v7z4-qsz8};
\href{https://arxiv.org/abs/2405.17557v2}{arXiv:2405.17557v2}.

\bibitem{convexroof}
C. Langrenez, S. De Bi\`evre and D. R. M. Arvidsson-Shukur,
\emph{Convex roofs witnessing Kirkwood--Dirac nonpositivity},
J. Phys. A: Math. Theor. \textbf{58}, 215301 (2025);
\href{https://doi.org/10.1088/1751-8121/add58f}{doi:10.1088/1751-8121/add58f};
\href{https://arxiv.org/abs/2407.04558}{arXiv:2407.04558}.

\bibitem{Cai}
L.-Y. Cai, Y.-H. Yang and Z.-J. Zheng,
\emph{Kirkwood--Dirac classical states based on discrete Fourier transform:
representation with directed graph},
\href{https://arxiv.org/abs/2603.13863v1}{arXiv:2603.13863v1} (2026).

\bibitem{Burkat}
J. Burkat and S. Strelchuk,
\emph{Classification of Kirkwood--Dirac positivity-preserving unitaries},
Phys. Rev. A \textbf{114}, 042446 (2026);
\href{https://doi.org/10.1103/rsy9-12fx}{doi:10.1103/rsy9-12fx};
\href{https://arxiv.org/abs/2502.11784v4}{arXiv:2502.11784v4}.

\bibitem{Spriet}
M. Spriet,
\emph{Characterizing the Kirkwood--Dirac positivity on second countable LCA groups},
J. Fourier Anal. Appl. \textbf{32}, 99 (2026);
\href{https://doi.org/10.1007/s00041-026-10309-4}{doi:10.1007/s00041-026-10309-4};
\href{https://arxiv.org/abs/2507.23628v3}{arXiv:2507.23628v3}.

\bibitem{Thio}
J. J. Thio, S. Yang, N. Yunger Halpern, S. De Bi\`evre, C. H. W. Barnes
and D. R. M. Arvidsson-Shukur,
\emph{Kirkwood--Dirac nonpositivity is a necessary resource for quantum computing},
Phys. Rev. Lett. \textbf{137}, 080601 (2026);
\href{https://doi.org/10.1103/x819-898d}{doi:10.1103/x819-898d};
\href{https://arxiv.org/abs/2506.08092v2}{arXiv:2506.08092v2}.

\bibitem{ThioEntropic}
J. J. Thio, S. De Bi\`evre, C. H. W. Barnes and D. R. M. Arvidsson-Shukur,
\emph{Tight quasiprobabilistic entropic uncertainty relations for discrete
quantum systems},
\href{https://arxiv.org/abs/2611.04009v1}{arXiv:2611.04009v1} (2026).

\bibitem{SLBD}
M. Spriet, C. Langrenez, R. Brummelhuis and S. De Bi\`evre,
\emph{What is special about the Kirkwood--Dirac distributions? Only they
produce natural conditional expectations},
\href{https://arxiv.org/abs/2511.01996v2}{arXiv:2511.01996v2} (2026);
journal version in Symmetry,
\href{https://doi.org/10.3390/sym18061008}{doi:10.3390/sym18061008}.
\end{thebibliography}
\endgroup
\end{document}
